%1st paragraph
%1
Table~\ref{accresults} shows the accuracy for all twelve combinations of features and clustering methods used.
%2
\deleted{Subjects A to AB served as the test dataset for four posture classes.}
\revised{Experiment participants A to AB served as the test dataset for four posture classes.}
\revised{Four posture classes yield a 25\% baseline accuracy under uniform random guessing.}
\revised{Accuracy exceeded the 25\% baseline in all cases in Table~\ref{accresults}.}
%3
\deleted{Depending on the feature and clustering combination, test accuracy was stable or varied across test subjects in Table~\ref{accresults}.}
\revised{Depending on the feature and clustering combination, test accuracy was stable or varied across test experiment participants in Table~\ref{accresults}.}
\deleted{Classification accuracy varied across test subjects because point-cloud density and body size differed from the three training subjects.}
\begin{deletedblock}
The accuracy is reported with the data from subjects A to AB serving as the test dataset.
Since we are classifying 4 posture classes, a baseline accuracy of 25\% or higher is expected. 
In all cases Table~\ref{accresults}, the accuracy exceeded the baseline.
\end{deletedblock}

%2nd paragraph
%4
Regarding the clustering methods, no clustering or DBSCAN tended to achieve higher accuracy than K-means.
  %4-1
  \revised{k-means tended to achieve lower accuracy than no clustering or DBSCAN for the same feature type.}
  \deleted{k-means tended to lower accuracy by removing hand-region points required to distinguish typing from pad operation.}

%3rd paragraph
%5
In terms of features, normals tended to yield higher accuracy when no clustering or DBSCAN was employed.
  %5-1
  \revised{Classification using normals yielded comparatively stable test accuracy in Table~\ref{accresults}.}
  \deleted{Normals yielded stable test accuracy when no clustering or DBSCAN was applied.}
  %5-2
  \deleted{Dimensionality features resulted in higher accuracies than the other three feature types for several test subjects outside the training set of subjects A, B, and C in Table~\ref{accresults}.}
  \revised{Dimensionality features resulted in higher accuracies than the other three feature types for several test experiment participants outside the training set of experiment participants A, B, and C in Table~\ref{accresults}.}
  \deleted{Dimensionality features achieved higher accuracy than other features for several subjects outside the training set when local shape resembled training examples.}

\begin{table*}[b]
  \centering
  \caption{Classification accuracy for each subject and feature and clustering combination.}
  \label{accresults}
  \footnotesize
  \setlength{\tabcolsep}{3pt}
  \resizebox{\textwidth}{!}{%
  \begin{tabular}{@{}c*{12}{c}@{}}
    \toprule
    & \multicolumn{3}{c}{Normals} & \multicolumn{3}{c}{Dim.\ features} & \multicolumn{3}{c}{Prop.\ variance} & \multicolumn{3}{c}{FPFH} \\
    \cmidrule(lr){2-4}\cmidrule(lr){5-7}\cmidrule(lr){8-10}\cmidrule(lr){11-13}
    Subject &
    \multicolumn{1}{c}{\label{normalnone}None} & \multicolumn{1}{c}{\label{normaldbscan}DBSCAN} & \multicolumn{1}{c}{\label{normalkmeans}k-means} &
    \multicolumn{1}{c}{\label{dimnone}None} & \multicolumn{1}{c}{\label{dimdbscan}DBSCAN} & \multicolumn{1}{c}{\label{dimkmeans}k-means} &
    \multicolumn{1}{c}{\label{proportionnone}None} & \multicolumn{1}{c}{\label{proportiondbscan}DBSCAN} & \multicolumn{1}{c}{\label{proportionkmeans}k-means} &
    \multicolumn{1}{c}{\label{fpfhnone}None} & \multicolumn{1}{c}{\label{fpfhdbscan}DBSCAN} & \multicolumn{1}{c}{\label{fpfhkmeans}k-means} \\
    \midrule
A & $0.65$ & $0.71$ & $0.60$ & $0.43$ & $0.86$ & $0.79$ & $0.70$ & $0.82$ & $0.70$ & $0.73$ & $0.85$ & $0.75$ \\
B & $0.98$ & $0.99$ & $0.88$ & $0.94$ & $0.97$ & $0.92$ & $0.98$ & $0.91$ & $0.82$ & $0.99$ & $0.96$ & $0.89$ \\
C & $0.70$ & $0.70$ & $0.70$ & $0.70$ & $0.70$ & $0.70$ & $0.70$ & $0.70$ & $0.70$ & $0.70$ & $0.70$ & $0.70$ \\
D & $0.61$ & $0.72$ & $0.41$ & $0.59$ & $0.67$ & $0.63$ & $0.47$ & $0.44$ & $0.63$ & $0.50$ & $0.60$ & $0.55$ \\
E & $0.63$ & $0.69$ & $0.46$ & $0.73$ & $0.60$ & $0.59$ & $0.60$ & $0.58$ & $0.56$ & $0.66$ & $0.53$ & $0.57$ \\
F & $0.72$ & $0.75$ & $0.54$ & $0.83$ & $0.58$ & $0.56$ & $0.61$ & $0.55$ & $0.62$ & $0.76$ & $0.58$ & $0.50$ \\
G & $0.68$ & $0.75$ & $0.40$ & $0.72$ & $0.53$ & $0.54$ & $0.56$ & $0.55$ & $0.54$ & $0.75$ & $0.55$ & $0.40$ \\
H & $0.71$ & $0.74$ & $0.68$ & $0.75$ & $0.68$ & $0.66$ & $0.60$ & $0.62$ & $0.61$ & $0.72$ & $0.63$ & $0.69$ \\
I & $0.76$ & $0.90$ & $0.63$ & $0.92$ & $0.66$ & $0.65$ & $0.61$ & $0.61$ & $0.60$ & $0.75$ & $0.64$ & $0.62$ \\
J & $0.72$ & $0.89$ & $0.88$ & $0.68$ & $0.90$ & $0.75$ & $0.73$ & $0.80$ & $0.84$ & $0.95$ & $0.81$ & $0.92$ \\
K & $0.74$ & $0.87$ & $0.87$ & $0.97$ & $0.86$ & $0.77$ & $0.65$ & $0.82$ & $0.79$ & $0.83$ & $0.80$ & $0.80$ \\
L & $0.73$ & $0.83$ & $0.90$ & $1.00$ & $0.86$ & $0.81$ & $0.63$ & $0.83$ & $0.80$ & $0.80$ & $0.81$ & $0.85$ \\
M & $0.59$ & $0.59$ & $0.45$ & $0.56$ & $0.60$ & $0.44$ & $0.56$ & $0.46$ & $0.40$ & $0.49$ & $0.60$ & $0.45$ \\
N & $0.67$ & $0.60$ & $0.62$ & $0.63$ & $0.76$ & $0.75$ & $0.84$ & $0.60$ & $0.68$ & $0.75$ & $0.84$ & $0.68$ \\
O & $0.85$ & $0.71$ & $0.71$ & $0.70$ & $0.76$ & $0.73$ & $0.73$ & $0.61$ & $0.68$ & $0.80$ & $0.73$ & $0.72$ \\
P & $0.79$ & $0.72$ & $0.70$ & $0.70$ & $0.82$ & $0.77$ & $0.68$ & $0.56$ & $0.75$ & $0.81$ & $0.80$ & $0.76$ \\
Q & $0.70$ & $0.67$ & $0.68$ & $0.70$ & $0.70$ & $0.64$ & $0.66$ & $0.60$ & $0.64$ & $0.71$ & $0.70$ & $0.68$ \\
R & $0.90$ & $0.88$ & $0.92$ & $0.72$ & $0.99$ & $0.84$ & $0.84$ & $0.89$ & $0.78$ & $0.80$ & $1.00$ & $0.87$ \\
S & $0.68$ & $0.72$ & $0.78$ & $0.92$ & $0.80$ & $0.73$ & $0.60$ & $0.79$ & $0.65$ & $0.59$ & $0.71$ & $0.83$ \\
T & $0.93$ & $0.87$ & $0.70$ & $0.70$ & $0.79$ & $0.67$ & $0.79$ & $0.70$ & $0.71$ & $0.85$ & $0.91$ & $0.74$ \\
U & $0.69$ & $0.81$ & $0.85$ & $0.71$ & $0.91$ & $0.76$ & $0.80$ & $0.89$ & $0.82$ & $0.87$ & $0.80$ & $0.86$ \\
V & $0.62$ & $0.62$ & $0.82$ & $0.74$ & $0.64$ & $0.66$ & $0.61$ & $0.57$ & $0.70$ & $0.54$ & $0.57$ & $0.66$ \\
W & $0.62$ & $0.63$ & $0.65$ & $0.32$ & $0.58$ & $0.56$ & $0.56$ & $0.57$ & $0.50$ & $0.41$ & $0.62$ & $0.41$ \\
X & $0.61$ & $0.65$ & $0.47$ & $0.36$ & $0.64$ & $0.48$ & $0.53$ & $0.60$ & $0.47$ & $0.41$ & $0.57$ & $0.43$ \\
Y & $0.60$ & $0.59$ & $0.64$ & $0.61$ & $0.67$ & $0.62$ & $0.60$ & $0.63$ & $0.60$ & $0.60$ & $0.65$ & $0.61$ \\
Z & $0.62$ & $0.59$ & $0.68$ & $0.77$ & $0.68$ & $0.58$ & $0.60$ & $0.64$ & $0.55$ & $0.66$ & $0.67$ & $0.66$ \\
AA & $0.64$ & $0.87$ & $0.94$ & $0.90$ & $0.79$ & $0.88$ & $0.62$ & $0.76$ & $0.71$ & $0.58$ & $0.54$ & $0.63$ \\
AB & $0.70$ & $0.90$ & $0.96$ & $0.82$ & $0.94$ & $0.87$ & $0.68$ & $0.94$ & $0.83$ & $0.87$ & $0.72$ & $0.86$ \\
    \bottomrule
  \end{tabular}%
  }
\end{table*}

%4th paragraph (C.2-1; after C.1-6 three-class paragraph when present)
\begin{revisedblock}
%6
We report four-class and three-class confusion counts on test data A.
%7
For normals without clustering, pad operation yielded 172 false predictions as mouse operation.
  %7-1
  Typing yielded 103 false predictions as pad operation.
%8
For normals with DBSCAN, pad operation yielded 30 false predictions as mouse operation.
  %8-1
  Typing yielded 230 false predictions as pad operation.
%9
For normals without clustering, pad-or-typing yielded 189 false predictions as mouse operation.
%10
For normals with DBSCAN, pad-or-typing yielded 30 false predictions as mouse operation.
%11
Full four-class and three-class confusion matrices for all twelve preprocessing configurations appear in the response letter.
\end{revisedblock}
